\section{Geschichte der Künstlichen Intelligenz}

Das Konzept intelligenter Maschinen existiert seit Jahrzehnten. Das Gebiet der Künstlichen
Intelligenz begann sich jedoch in den 1950er Jahren formal zu konsolidieren, insbesondere nach den
Arbeiten von Alan Turing und der Dartmouth-Konferenz im Jahr 1956.

In den folgenden Jahrzehnten suchten Forscher nach der Entwicklung von Algorithmen, die in der Lage waren,
logisches Denken und Entscheidungsfindung zu simulieren. Anfangs basierten viele Systeme auf
festen Regeln, die manuell von Spezialisten definiert wurden.

In den 1990er Jahren ermöglichte die Zunahme der Rechenleistung erhebliche Fortschritte beim
statistischen Lernen. Später befeuerte das Wachstum des auf dem Internet verfügbaren Datenvolumens
die Entwicklung von Deep-Learning-Techniken noch weiter.

Aktuell werden Sprachmodelle und tiefe neuronale Netze in Aufgaben der Computer Vision,
der natürlichen Sprachverarbeitung und der prädiktiven Analyse eingesetzt.
 